%% ─────────────────────────────────────────────────────────────────────────────
%%  Variación acotada clásica — caso real
%% ─────────────────────────────────────────────────────────────────────────────
\section{Variación acotada clásica}

% IDEA ORAL: Comenzar recordando que todo el mundo conoce la variación
% acotada en el caso real. Esto es el punto de partida; la pregunta
% que motiva el trabajo es qué ocurre cuando el codominio no es ℝ.

\begin{frame}{Particiones y variación respecto a una partición}
\framesubtitle{Caso clásico: $f:[a,b]\to\mathbb{R}$}

Una \textbf{partición} de $[a,b]$ es un conjunto ordenado de puntos
\[
\mathcal{P}\colon \{ a = x_0 < x_1 < \cdots < x_n = b \}.
\]

\medskip

Dada $f:[a,b]\to \mathbb{R}$, la \textbf{variación de $f$ respecto a $\mathcal{P}$} es
\[
V(f,\mathcal{P})=\sum_{k=1}^n \bigl|f(x_k)-f(x_{k-1})\bigr|.
\]

\medskip

\begin{alertblock}{Idea intuitiva}
$V(f,\mathcal{P})$ suma los desplazamientos verticales de $f$ sobre cada
subintervalo de la partición.
\end{alertblock}

\end{frame}

%──────────────────────────────────────────────────────────────────────────────

\begin{frame}{Variación total y variación acotada}
\framesubtitle{Definición central}

\begin{definition}[Variación total]
La \textbf{variación total} de $f$ en $[a,b]$ es
\[
V_a^b(f)=\sup\bigl\{ V(f,\mathcal{P}) : \mathcal{P} \text{ partición de } [a,b]\bigr\}.
\]
Si $V_a^b(f)<\infty$, $f$ es de \textbf{variación acotada} y escribimos
$f\in BV[a,b]$.
\end{definition}

\end{frame}

%──────────────────────────────────────────────────────────────────────────────

\begin{frame}{Ejemplos fundamentales}
\framesubtitle{Qué funciones reales tienen variación acotada}

\begin{enumerate}
    \item \textbf{Función constante.}
    $f(x)\equiv c \;\Longrightarrow\; V_a^b(f)=0$.

    \item \textbf{Función monótona.}
    Si $f$ es creciente, $V(f,\mathcal{P})=f(b)-f(a)$ para toda partición,
    por lo tanto $V_a^b(f)=f(b)-f(a)<\infty$.

    \item \textbf{Función Lipschitz.}
    Si $|f(x)-f(y)|\le L|x-y|$, entonces $V_a^b(f)\le L(b-a)<\infty$.
\end{enumerate}

% IDEA ORAL: Resaltar que toda función monótona acotada es de v.a.
% Este hecho se pierde en el caso vectorial, y esa es exactamente
% la pregunta central del trabajo.

\end{frame}

%──────────────────────────────────────────────────────────────────────────────

\begin{frame}{Propiedades de la variación acotada}
\framesubtitle{Estructura algebraica de $BV[a,b]$}

\begin{block}{Lema}
Sean $f,g \in BV[a,b]$, $c\in\mathbb{R}$. Entonces:
\begin{enumerate}
    \item $V_a^b(f) = 0 \iff f$ es constante.
    \item $V_a^b(cf) = |c|\,V_a^b(f)$.
    \item $V_a^b(f+g) \le V_a^b(f) + V_a^b(g)$.
    \item $V_a^b(fg) \le \|f\|_{\infty} V_a^b(g) + \|g\|_{\infty} V_a^b(f)$.
    \item $V_a^b(|f|) = V_a^b(f)$.
    \item Si $c\in(a,b)$, entonces $V_a^b(f)=V_a^c(f)+V_c^b(f)$.
\end{enumerate}
\end{block}

\end{frame}

%──────────────────────────────────────────────────────────────────────────────

\begin{frame}{$BV[a,b]$ como espacio de Banach}
\framesubtitle{Completitud y norma canónica}

\begin{theorem}
$BV[a,b]$ es un espacio de Banach con la norma
$\|f\|_{BV}= |f(a)| + V_a^b(f)$.
\end{theorem}

\end{frame}

%──────────────────────────────────────────────────────────────────────────────

\begin{frame}[shrink=5]{Teorema de descomposición de Jordan}
\framesubtitle{Resultado fundamental del caso real}

\begin{theorem}[Jordan, 1881]
Una función $f:[a,b]\to\mathbb{R}$ es de variación acotada \emph{si y solo si}
puede escribirse como
\[
f = g - h,
\]
donde $g$ y $h$ son funciones \textbf{crecientes} en $[a,b]$.
\end{theorem}

\smallskip

\begin{alertblock}{Pregunta central de esta investigación}
¿Cuándo se puede garantizar una descomposición análoga cuando
$f:[a,b]\to X$, con $(X,\|\cdot\|,\Xp)$ un espacio normado ordenado?
La respuesta depende de la \textbf{geometría del cono} $\Xp$.
\end{alertblock}

% IDEA ORAL: Este es el punto de giro de la presentación. Todo lo que
% sigue busca responder precisamente esta pregunta.

\end{frame}

%──────────────────────────────────────────────────────────────────────────────

\begin{frame}{¿Qué cambia al pasar a funciones vectoriales?}
\framesubtitle{Las dificultades del caso vectorial}

En $(\mathbb{R},|\cdot|,\mathbb{R}_{+})$:
\begin{itemize}
  \item[\checkmark] Toda función monótona es de variación acotada.
  \item[\checkmark] Toda función de v.a.\ admite descomposición de Jordan.
  \item[\checkmark] Solo hay \emph{una} noción de variación acotada.
\end{itemize}

\medskip

Al reemplazar $\mathbb{R}$ por $(X,\|\cdot\|,\Xp)$:
\begin{itemize}
  \item[$\times$] \textbf{Tres nociones distintas} de variación acotada.
  \item[$\times$] Las funciones crecientes \emph{no necesariamente} tienen v.a.
  \item[$\times$] La descomposición de Jordan \emph{no siempre} existe.
\end{itemize}

\smallskip

\textbf{Determinante:} las propiedades del cono $\Xp$
(normalidad, solidez, ser generador, ser extendible).

\end{frame}
